% 本节围绕幻觉抑制、受约束执行和Context Ledger恢复组织机制。
\section{神经符号协同的幻觉抑制与安全执行}
\label{sec:mechanism_coordination}

\subsection{安全目标与分层校验架构}

长程协作中的幻觉风险不仅表现为答案中的错误陈述，还包括引用不存在的资源、将未经证实的阶段结果写入共享状态，以及把执行异常解释为任务成功。本节围绕\textbf{候选生成、证据核验、受约束执行与状态修复}组织神经符号协同，使模型提出的建议经过明确检查后，才能影响后续任务和最终交付。

本文所讨论的安全性包括证据使用的可靠性、已编码约束下的执行合规性和错误状态的传播控制。事实正确、来源支持、计划可行与权限合法分别检查；某一层通过不替代其他层的判断。系统以降低无依据结论被采纳以及不合规动作被执行的风险为目标，具体降低幅度由对照实验确定。

\begin{table}[htbp]
\centering
\caption{幻觉与执行风险的分层处理}
\label{tab:safety_risks}
\begin{tabular}{p{0.20\textwidth}p{0.36\textwidth}p{0.33\textwidth}}
\toprule
风险 & 检查机制 & 处理结果 \\
\midrule
计划或资源虚构 & PlanIR结构校验、资源目录和权限过滤 & 返回反例，修订或暂停 \\
结论缺少证据 & 主张分解、来源核验及外部验收记录 & 补证、标记未知或拒绝提交 \\
动作突破边界 & 调用前检查、可信域及预算可行域 & 阻断、改选或请求授权 \\
错误跨节点传播 & 候选结果隔离、保护通知和版本检查 & 禁止错误进入可复用状态与后继输入 \\
静默失效与中断 & 任务进展检测、Context Ledger与Checkpoint恢复 & 核对状态后有限修复或退出 \\
\bottomrule
\end{tabular}
\end{table}

TIDE-Anchor提供当前有效目标、状态和证据入口，TIDE-Mesh控制协作边及消息内容，TIDE-Placement在可行资源范围内安排推理。本节的校验与恢复流程连接三者，分别决定候选是否可以执行、产物是否可以采纳，以及失效后哪些状态仍可继续使用。

\begin{figure}[htbp]
\centering
\fbox{\parbox{0.92\textwidth}{\centering
模型生成候选计划与结果\\[3pt]
$\downarrow$ PlanIR约束检查 $\downarrow$\\[3pt]
证据核验 $\rightarrow$ 动作与可信域检查 $\rightarrow$ 受控执行\\[3pt]
$\downarrow$ 产物验收、状态版本检查 $\downarrow$\\[3pt]
提交Context Ledger $\rightarrow$ 定向通信与经验准入\\[3pt]
失败或争议 $\rightarrow$ 隔离、补证、有限修订或一致恢复}}
\caption{以幻觉抑制与安全执行为目标的协同校验流程}
\label{fig:safety_validation_loop}
\end{figure}

\subsection{PlanIR约束校验与反例驱动修订}
\label{sec:neuro_symbolic_coordination}

计划校验将模型的自然语言建议转化为可检查的执行条件。PlanIR记录目标、子任务、依赖、约束、资源、验收要求与版本。模型生成候选后，SymbolicPlanValidator先检查标识唯一性、依赖存在性、自依赖、负时长与依赖环，避免资源调度和团队执行消费结构无效的计划。

确定性后端进一步依据最早调度检查已编码的角色、单任务资源、先后、互斥、截止时间和必需证据条件。可选Z3后端对已支持的时序约束搜索可行解并返回反例。不同后端实际覆盖的检查集合和调度语义分别记录；PlanIR中的描述性前置条件、效果和软约束不自动获得证明。尤其是计划中存在“测试证据”这一验收条目，只表示计划提出了证据要求，不能证明测试已经执行并通过。

校验失败后，将约束标识、违反原因和见证信息反馈给修订器，在有限次数内重新提交。新计划或修订版本进入执行前再次校验，旧版本的通过记录不能直接授权新版本继续。运行状态绑定计划摘要与版本；版本变化后重新确认任务完成状态，避免沿用已失效的执行许可。

\begin{algorithm}[htbp]
\caption{反例驱动的有限计划校验与修订}
\label{alg:neuro_symbolic_validation}
\KwIn{目标$g$，上下文$c$，生成器$F$，校验器$V$，修订器$R$，修订上限$J$}
\KwOut{计划$p$，校验结果$v$及修订轨迹}
$p\leftarrow F(g,c)$\;
\For{$j\leftarrow0$ \KwTo $J$}{
    $v\leftarrow V(p)$\;
    \If{$v$通过}{记录计划版本及摘要；\Return{$p,v$与修订轨迹}\;}
    将反例写入保护反馈，保留约束标识与见证\;
    \If{$j=J$}{\textbf{break}\;}
    $p\leftarrow R(p,v.\mathrm{counterexamples})$，更新版本\;
}
记录计划校验暂停状态及未解决原因\;
\end{algorithm}

\subsection{证据驱动的事实核验与产物准入}

事实核验以可定位主张及其证据关系为对象。FActScore按原子事实检查来源支持，RARR通过检索和修订处理缺少依据的内容\cite{tide_factscore,tide_rarr}。本文借鉴这一组织方式，将关键数值、日期、引用、测试结论和执行完成声明拆为可核验项，并关联原始文档、工具调用或外部验证记录。

每个主张保留“有证据支持”“与证据冲突”或“信息不足”的判定。来源可访问仅表示证据入口有效，仍需检查证据是否支持该主张、是否适用于当前时间与任务版本。多个Agent复述同一段文本不增加独立证据数量；模型自报置信度和HTTP成功也不能替代事实核验。

结构化验收使用必需字段、产物存在性、测试结果和外部证据标识等规则；语义支持关系可由评审模型辅助判断，并对重要结论使用独立来源或任务专用验证器。外部验证记录从执行前的受控状态读取，节点不能通过在自身输出中声明“已验证”来批准自己。证据不足时先补证或保留未知状态，不能为了获得完整答案而补造结论。

阶段验收分数沿用第\ref{sec:tide_mesh}节的规则与模型组合，但硬性证据缺失、版本失效或明确冲突具有否决作用。通过后，产物才更新Context Ledger中的可用任务结果，并按第\ref{sec:tide_anchor}节规则提炼可复用记忆；未通过产物保留为候选或失败记录，阻止其进入正常的记忆准入和胶囊发送路径。未配置事实验证器的通用输出检查，只能表明通过已有规则，不称为事实真实性证明。

\subsection{调用前约束与跨节点传播控制}

执行安全在工具或模型调用前生效。控制器读取当前有效的授权、资源、输入条件和预算，检查目标工具及参数是否满足已编码规则。资源目录限定可调用对象；Placement的可行域排除不可信、未授权、能力不符或预算不足的候选。学习收益和质量分不能抵消这些禁止条件，策略故障回退也在同一可行域内选取。

对涉及外部数据处理的任务，安全预处理和实际请求载荷需对应，调用前再次核对策略与资源状态。网页、检索材料和工具输出属于业务证据，不能凭其中包含的自然语言指令改变用户授权和系统边界。可编码的字段由显式检查处理，未编码的语义风险不能仅凭形式校验排除。

传播控制由验收隔离与通信协议共同承担。获准传递的产物携带来源、验证状态、任务版本及必要证据引用；压缩不改变主张、约束变化和关键来源。撤销、失败与冲突属于保护信息，接收者按已确认版本校验状态增量。Context Ledger记录的事实被撤销后，后续上下文不再将其列为已确认依据；已消费该事实的任务应进入重新核验或修订路径。

能力变化、证据争议和阶段返工经任务控制器更新剩余协作关系，服务失败则经Placement更新资源状态。三者共享失败原因与版本信息，避免同一错误在角色切换或资源回退后重新被当作正确前提。

\subsection{Context Ledger驱动的状态恢复与局部修复}

恢复以Context Ledger维护的任务状态为中心，以Checkpoint保存的一致版本为载体。Context Ledger描述原始目标、当前约束、进度、事实验证状态与失败记录；Checkpoint同时记录该视图、执行前沿及必要消息状态。恢复时需要重建任务信息与继续执行的位置，而不仅是重新加载一段历史摘要。

统一恢复入口读取同一图检查点中的Context Ledger视图与通信状态，先核对版本，再协调当前有效约束和已撤销事实。可表示为
\begin{equation}
\widehat{S}=\operatorname{Reconcile}
\left(\operatorname{Load}(K_j),\mathcal{C}_{\mathrm{current}},
\mathcal{R}_{\mathrm{revoked}},\mathcal{U}_{\mathrm{external}}\right),
\label{eq:safety_reconcile}
\end{equation}
其中$K_j$为恢复点，$\mathcal{C}_{\mathrm{current}}$为最新有效约束，$\mathcal{R}_{\mathrm{revoked}}$为事实撤销信息，$\mathcal{U}_{\mathrm{external}}$为外部执行结果。该协调过程按对象版本和权限处理冲突，不以恢复旧版本为由撤销新的用户限制。

错误修复分为补证、当前节点修订、依赖范围内重执行和中断恢复。若A提取数据、B计算、C写作，C发现B使用了失效单位，应先核对原文和事实版本；确认B错误后保留A的有效产物，修订B并重新生成受影响的C。只有错误来自A或其证据时，才扩展重执行范围。当前运行时通过既有暂停、修订和重路由接口执行明确的修复目标；完整的自动依赖闭包分析属于进一步扩展。

可重试服务故障只需要更新资源监控并重新申请；成员能力不足则重新选人；阶段质量不达标则补证或返工。以上情况不必都恢复整个Checkpoint。对于持续返回正常状态但缺少新产物的情况，进展检测触发核验，并结合长计算心跳区分正常等待与静默停滞。

具有外部副作用的步骤还需核对执行记录和幂等标识。恢复Context Ledger不能撤销已经发送的通知、提交的订单或外部写入；需要补偿或人工处理时，记录原因并暂停。所有修订、重试、转交与恢复均设置预算，恢复后必须重新通过当前验收和约束检查才计为成功。

\subsection{闭环算法与评价口径}

\begin{algorithm}[htbp]
\caption{幻觉抑制与安全执行闭环}
\label{alg:safety_execution_loop}
\KwIn{共享状态$S$，计划$p$，当前节点$n$及有限修复预算$b$}
\KwOut{已验收状态或带原因的暂停/失败状态}
核对计划版本，必要时重新执行PlanIR校验\;
从Context Ledger读取有效目标、约束与证据状态\;
构造节点上下文，检查动作及资源可行域\;
执行节点，保留候选结果与实际调用记录\;
核验产物契约、必需证据及有效版本\;
\eIf{验收通过}{
    更新Context Ledger；按作用域总结记忆并发送受控胶囊\;
    保存状态、执行位置与消息接收基准\;
}{
    隔离候选产物，保留冲突和失败原因\;
    在$b$内补证、修订、重路由或恢复，否则暂停\;
}
复核实际任务进展与剩余预算，返回闭环状态\;
\end{algorithm}

评价同时观察结果的依据、执行边界和有效交付。不能仅凭拒绝更多问题来宣称幻觉减少，也不能以记录了恢复事件替代恢复成功。各项口径见表~\ref{tab:mechanism_observation}。

\begin{table}[htbp]
\centering
\caption{幻觉抑制与安全执行的运行证据及评价内容}
\label{tab:mechanism_observation}
\begin{tabular}{p{0.20\textwidth}p{0.36\textwidth}p{0.33\textwidth}}
\toprule
层次 & 运行证据 & 评价内容 \\
\midrule
计划约束 & 计划版本、后端、违反项与反例 & 已编码违规检出、误拒绝及修订结果 \\
事实依据 & 主张、来源、核验结果与未知项 & 证据支持比例、无依据结论采纳及任务覆盖 \\
执行安全 & 允许动作、阻断原因与实际调用 & 硬约束违规、有效调用与错误放行 \\
传播控制 & 验收状态、消息版本及撤销记录 & 错误下游采纳、版本一致与必要信息保留 \\
恢复与交付 & Context Ledger版本、执行前沿及最终产物 & 恢复成功、重复副作用、交付质量及总成本 \\
\bottomrule
\end{tabular}
\end{table}

对照实验分别移除计划校验、证据准入、传播隔离或Context Ledger恢复，在相同任务与故障条件下比较结果。测试包括虚构引用、未执行却声称通过的测试、失效事实、冲突消息、计划替换、资源限流与中断续跑。自动核验需要以人工抽样或任务标准校准，并同时报告误报、漏报、任务完成和新增时延。

现有工程已接入PlanIR校验、反例修订、规则验收、Context Ledger、有限恢复及资源安全过滤。本轮实现补齐计划变更重验、验收失败隔离、外部证据标识检查、Context Ledger恢复协调与通信版本恢复等接口。通用自然语言事实核验、跨工具的自动依赖修复和广泛领域中的幻觉降低幅度仍需专门实现和实验；形式可行、证据存在与业务正确分别评价。

